\subsubsection{Generación y reserva efectiva de combustible (GEN-4)}
\label{sec:sd4-generacion-gen4}

GEN-4 desarrolla RT-06.08 desde el parque completo y las cargas de recuperación. Cada FG Wilson \texttt{P200-6/FGL30120} atiende exclusivamente su camino; se conserva PRP de 180 kVA/144 kW a 400/230 V, 50 Hz. No se emplea la potencia standby para justificar veinticuatro horas de servicio. La selección ambiental de proyecto fija admisión máxima de 35\,$^{\circ}$C, altura de hasta 100 m, aire exterior sin recirculación de radiador y escape independiente. Estos son límites de instalación nueva; no condiciones observadas ni capacidad corregida del fabricante. La ficha nominal usa 25\,$^{\circ}$C, 100 m y 30\,\% HR. Por tanto, el radiador anunciado para 50\,$^{\circ}$C no demuestra ausencia de reducción de potencia a la admisión de proyecto \parencite[pp.~1--5]{p7-fg-p200-ficha}.

La demanda permitida por grupo, incluidos rectificador, su recarga, tratamiento ambiental asignado y auxiliares, conserva los dos límites simultáneos de 90 kW y 90 kVA. La potencia corregida requerida es al menos 129 kW y 90 kVA. En consecuencia, la corrección activa admisible respecto de 144 kW nominales es como máximo $1-129/144=10,4167\,\%$; una corrección mayor obliga a sustituir el grupo y a recalcular consumo, estanque y obra. El criterio propio de media a lo sumo 70\,\% de PRP efectivo conserva $90/129=69,7674\,\%$. Se exige además comprobación de capacidad reactiva, desequilibrio y armónicos; esos cocientes activos no acreditan excitación ni respuesta a escalones.

La entrada del rectificador a 400 V se limita a la cota publicada de 108 A, equivalente a $\sqrt{3}(400)(108)/1000=74,8246$ kVA. Este valor contiene el cargador y no admite sumar de nuevo 11,76 kW de recarga. La envolvente adicional externa a esa entrada es $90-74,8246=15,1754$ kVA por grupo, incluyendo auxiliares y cualquier climatización que se conecte antes de UPS. Una carga ambiental alimentada después de UPS debe caber además en sus 60 kVA/60 kW y en cada fase. El tratamiento de bancos no puede asignarse fuera de la UPS para mantenerlo disponible durante el intervalo previo al arranque sin fuente adicional. Antes de liberar compra, se registra para cada auxiliar fuente durante red, batería, grupo y recarga; se recalculan ambas fronteras de manera separada \parencites[pp.~33--35]{p7-eaton-93t-guia}{p7-eaton-93t-tecnica}.

La gestión autónoma arranca ambos grupos al perder la red, confirma tensión y frecuencia y transfiere cada entrada por separado. La UPS mantiene las cargas mientras califica la fuente; no se toma el cierre del contactor como recuperación de autonomía. El servicio autorizado configura un límite de entrada y una rampa de rectificador compatibles con el grupo seleccionado. Se secuencian arranque climático, recuperación de banco y maniobras de barrera, conservando las prioridades de seguridad. Los parámetros de la rampa, el reparto efectivo entre salida y recarga hasta 60 A DC y la autorización de inhibición de recarga por ventilación siguen por individualizar: no se copia una rampa de otra familia UPS.

Para evitar operación prolongada a baja carga, GEN-4 añade una carga resistiva exterior por grupo, independiente de la bancada de aceptación de UPS. Se prescribe capacidad de 60 kW con escalones de hasta 5 kW y desconexión prioritaria antes de una nueva demanda de servicio. Como regla propia preliminar, el control persigue al menos $0,40(144)=57,6$ kW en generación estable y mantiene techo de 90 kW/90 kVA incluyendo la bancada, nunca 90 kW de servicio más 60 kW de bancada. La selección del control registra potencia y umbral, evita reconexión cíclica y no retira cargas críticas para cumplir el mínimo. El fabricante debe confirmar el régimen admisible y si requiere un umbral o ciclo diferente; la regla del cuarenta por ciento no se presenta como instrucción oficial ni como solución ya recibida. GCI-4 selecciona Avtron 4100/60 kW/400 V/50 Hz con ALC; quedan por acreditar motor y control auxiliar a 50 Hz, escalones, CT, protecciones y régimen autorizado del grupo.

El combustible inicial no depende de reposición en las primeras veinticuatro horas. GCI-4 amplía cada depósito a 2.000 L nominales y al menos 1.500 L útiles: 200 L de expansión y hasta 300 L no extraíbles. A+B totaliza 3.000 L útiles sin compartir reserva para acreditar un camino. A 41,3 L/h publicados resultan 36,3196 h; con margen propio del veinte por ciento, el consumo corregido admisible es 1.500/(24$\times$1,20)=52,0833 L/h. Un consumo superior obliga a ampliar el depósito antes de compra. La bancada participa del consumo y los estados del grupo; mayor volumen no demuestra derating ni respuesta transitoria \parencite{p7-fg-p200-ficha}.

Se conserva Enex Express declarado, con Copec Empresas alternativo, confirmación en una hora y entrega en ocho horas desde pedido inicial. El nuevo umbral es 600 L útiles; a 52,0833 L/h quedan 183,3333 L/3,52 h al vencimiento, y a 41,3 L/h quedan 269,6 L/6,5278 h. La entrega restituye 1.500 L útiles por camino. El acuerdo se formaliza antes de habilitación; no se afirma contrato firmado. RT-06.08 sigue por completar por selección ambiental/transitoria del grupo, control auxiliar de bancada y demanda completa de bancos. GCI-4 desarrolla la reserva vigente; ENE-4 conserva sus escenarios previos como referencia condicionada.

\subsubsection{Instalación independiente y memoria de distribución (ELE-4)}
\label{sec:sd4-electrica-ele4}

ELE-4 fija una instalación nueva de 400/230 V, 50 Hz, tres fases, neutro y PE, con empalme y tableros exclusivos del recinto. Ningún circuito se toma del tablero de oficina ni se comparten sus neutros. La independencia funcional no elimina la equipotencialidad de las masas accesibles: los enlaces de protección se diseñan para evitar diferencias peligrosas con el resto de la instalación. El esquema elegido aguas abajo de las fuentes es TN-S, con PE separado de N; la transferencia omnipolar selecciona fases y neutro sin interrumpir PE. La conexión de servicio de cada fuente y la continuidad del neutro a través de UPS deberán concordar con el esquema autorizado Eaton antes de fabricar tableros \parencites{lote10-sec-ric}{lote13-sec-ric05}{lote13-sec-ric06}[pp.~33--35]{p7-eaton-93t-guia}.

El unifilar de adquisición establece una barra de red exclusiva, protección general omnipolar de 320 A y dos salidas independientes de 160 A hacia A/B. Cada salida entra al puerto red de una transferencia automática de cuatro polos y al menos 320 A; el puerto grupo recibe su P200-6 mediante protección limitada a 160 A. No se paralelizan red/grupo ni A/B. Después de cada transferencia se ubican medición, seccionamiento, protección de rectificador de 125 A y UPS 60 kW; su salida pasa por seccionamiento/protección de 100 A al tablero propio A o B. El bypass tiene circuito y seccionador independiente abierto y bloqueado en operación Frequency Converter; no forma un puente entre A/B. Los auxiliares previos a UPS tienen circuitos exclusivos identificados y no adquieren continuidad en batería por estar dentro del módulo. Los bancos DC conservan protección por cadena y conjunto, neutro central y la interfaz TB20; no se conectan al PE como si fueran neutro AC.

\begin{table}[htbp]
\centering\small
\caption{Circuitos troncales de ELE-4 por camino}\label{tab:sd4-ele4-troncales}
\begin{tabularx}{\linewidth}{X R{16mm} L{31mm} R{16mm}}
\hline
\EncabezadoTabla{Circuito} & \EncabezadoTabla{Longitud máxima} & \EncabezadoTabla{Cu: fases / N / PE} & \EncabezadoTabla{Protección} \\ \hline
Nuevo empalme a barra exclusiva & 20 m & 150 / 300 / 150 mm$^2$ & 320 A \\ \hline
Barra exclusiva a transferencia; grupo a transferencia & 25 m cada uno & 70 / 120 / 70 mm$^2$ & 160 A \\ \hline
Transferencia a rectificador UPS & 20 m & 35 / 70 / 35 mm$^2$ & 125 A \\ \hline
UPS a tablero A o B & 15 m & 35 / 70 / 35 mm$^2$ & 100 A \\ \hline
Bypass separado, bloqueado & 20 m & 35 / 70 / 35 mm$^2$ & 100 A \\ \hline
\end{tabularx}
\FuenteTabla{Condiciones de obra de Synaptix. Conductores unipolares de tendido fijo RZ1-K, 0,6/1 kV, sin armadura, libres de halógenos, aislación de 90\,$^{\circ}$C; método F en bandeja perforada, sin sol directo, ambiente máximo de 30\,$^{\circ}$C y rutas A/B separadas. No corresponden a longitudes medidas ni cables ya comprados.}
\end{table}

El empalme exclusivo se reserva para 320 A, equivalente a 221,70 kVA a 400 V, frente al techo conjunto de fuentes de $2(90)=180$ kVA. El troncal común emplea 150 mm$^2$ de fase/PE y 300 mm$^2$ de neutro; en método F, la tabla 4.4 da 464/736 A, y la reserva combinada propia de 0,80 deja 371,2/588,8 A. Por tanto, la protección de 320 A no exige operar el conductor de fase sobre esa cota. La terminación se realiza en barra de tablero preparada para esas secciones, no en bornes UPS. La derivación de 160 A limita intencionadamente el uso de la placa del grupo; no se dimensiona su cable como si se permitieran los 259,81 A nominales completos.

RIC 04, tabla 4.4 para 90\,$^{\circ}$C y método F, publica 176 A para 35 mm$^2$, 279 A para 70 mm$^2$ y 400 A para 120 mm$^2$ con tres conductores cargados. Para presupuesto se impone un factor propio combinado de 0,80 como reserva frente a armónicos/neutro y disposición: resultan 140,8/223,2/320 A. Por tanto, $108\leq125\leq140,8$ A en entrada UPS, $86\leq100\leq140,8$ A en salida y $160\leq223,2$ A antes de transferencia. La reserva propia no reemplaza los factores normativos de la instalación real; si su producto es inferior a 0,80 se cambia disposición/sección antes de ejecutar. Los neutros de 70 mm$^2$ superan $1,7(35)=59,5$ mm$^2$ de la guía Eaton y caben en su límite AC de 70 mm$^2$; el PE conserva sección igual a fase. La guía recomienda 35 mm$^2$ para 60 kW y admite hasta 70 mm$^2$ en AC; no se intenta insertar 95/120 mm$^2$ en el borne UPS. La temperatura admisible de terminaciones y la solicitación del neutro se confirman con la variante suministrada; un cable de 90\,$^{\circ}$C no autoriza por sí solo el borne a esa temperatura \parencites[tabla~4.4, pp.~19--21]{lote10-sec-ric04}[pp.~34--35]{p7-eaton-93t-guia}.

El cálculo de caída emplea resistividad propia conservadora $\rho=0,0225$ $\Omega$mm$^2$/m y reactancia presupuestada de 0,08 $\Omega$/km. En la entrada UPS, aun tomando coseno y seno iguales a uno como cota, $\Delta V\leq\sqrt{3}(108)[0,0225(20)/35+0,00008(20)]=2,7044$ V, o 0,6761\,\% de 400 V. El tramo anterior de 25 m/70 mm$^2$ a 108 A añade 0,4693\,\% como máximo. El troncal común a 320 A añade 0,6374\,\%; la suma conservadora hasta rectificador es 1,7828\,\%. La salida de 15 m, a 86 A monofásicos con retorno de igual sección de fase como peor caso resistivo, añade $2(86)(0,0225)(15)/35=1,6586$ V, o 0,7211\,\% de 230 V. La cota resistiva de neutro real de 70 mm$^2$ es menor; el estudio final conserva reactancia/desbalance y armónicos. Los presupuestos de entrada y salida se separan por la regulación del inversor; no se suman a través de la conversión. Se reserva como objetivo propio 2\,\% hasta tablero de carga y 1\,\% adicional en circuitos finales; una ruta exterior larga requiere cálculo independiente, no prolongación automática de estos tramos.

Se fija una malla de proyecto bajo el módulo de 10 por 7 m: anillo de 34 m y dos conductores transversales de 10 y 7 m, total 51 m de cobre desnudo de 70 mm$^2$, enterrados a 0,8 m; cuatro electrodos de 3 m con registros en esquinas y barra principal de tierra accesible. La adquisición incluye enlaces equipotenciales de 35 mm$^2$ a tableros, UPS, gabinetes y estructuras, y de 70 mm$^2$ a generadores y malla. Se adopta resistencia objetivo de proyecto a lo sumo 5 $\Omega$, sin tratarla como comprobación de tensiones de paso/contacto o tiempo de despeje. Se mide resistividad del terreno nuevo para validar electrodo y corrosión; se amplía la malla antes de habilitación si falla el objetivo o la seguridad calculada. PE no usa contactos de transferencia ni separa tierras de A/B en islas. No se establece otra unión N--PE aguas abajo de las uniones de servicio aprobadas \parencite[secc.~5--8 y 12]{lote13-sec-ric06}.

El presupuesto de cortocircuito fija hasta 10 kA simétricos en barras de entrada y exige poder de corte al menos 25 kA a 400 V para protecciones troncales; el límite se recibe del nuevo empalme y del generador, y debe recalcularse si se excede. Para cobre XLPE se usa $k=143$ A$\sqrt{\mathrm{s}}$/mm$^2$ como hipótesis adiabática de selección. El menor conductor troncal de 35 mm$^2$ tolera presupuesto $k^2S^2=25.050.025$ A$^2$s; 10 kA durante 0,20 s implican 20.000.000 A$^2$s. La adquisición exige energía pasante inferior a ese límite en todo el rango y comprueba también PE y neutro. La secuencia nominal 160/125/100 A no acredita selectividad. En batería, el inversor limita corriente y no asegura disparo magnético de una protección de red; deben individualizarse corriente/tiempo de cortocircuito UPS y curvas de los circuitos finales. Las variantes, ajustes y tabla de selectividad entre cabecera, UPS, clima, PDU, borde y distribución final siguen por seleccionar. Los diferenciales se eligen por forma de fuga y condición de UPS, con protección personal en circuitos finales exigibles y discriminación aguas arriba; no se impone un único tipo AC a todo el parque.

\textbf{Protecciones troncales elegidas (ELE-19).} Se seleccionan conjuntos Schneider Electric ComPacT NSX de cuatro polos, bastidor F de 36 kA a 380/415 V y MicroLogic 5 E, adquiridos como bastidor más unidad de disparo compatible. La cabecera usa un \texttt{C40F4} con \texttt{C4045E400} (5.3 E/400 A), ajuste largo 320 A. Las dos derivaciones de red y las dos de grupo usan cuatro \texttt{C25F4} con \texttt{C2545E250} (5.2 E/250 A), ajuste largo 160 A. Las dos entradas UPS, dos salidas y dos bypass bloqueados usan seis \texttt{C16F4} con \texttt{C1645E160} (5.2 E/160 A), ajustes largos de 125 A en entrada y 100 A en salida/bypass. Son once conjuntos instalados; se dispone además de un conjunto de cada tamaño apagado en reserva, catorce adquiridos. Cada conjunto conserva interruptor y unidad de disparo como partes distintas del inventario, sin contar dos interruptores por conjunto. El catálogo publica los bastidores y unidades exactos, y ajuste de MicroLogic de 0,4 a 1 de In con ajuste fino de 1 A. El poder de corte supera la exigencia de 25 kA de ELE-4 en red; no sustituye el análisis de corriente mínima de falla \parencite[pp.~43, 49, 289 y 316]{lote19-schneider-compact}.

Los conjuntos se reciben con protección del neutro compatible con el neutro reforzado y con terminales admitidos para cada sección, a temperatura de servicio. Las funciones LSI permiten ajustar tiempo largo, corto e instantáneo; la lista de aparatos no acredita su coordinación. En particular, $160/125=1,28$ y $125/100=1,25$ son inferiores a la razón de 1,3 usada por el fabricante como regla usual de selectividad larga entre unidades electrónicas. Para selectividad energética tampoco se alcanza automáticamente la razón de bastidores 2,5: 400/250=1,60 y 250/160=1,5625. Se exige tabla/curvas de esas parejas con tolerancias, energía pasante, retardo y corriente disponible de cada fuente; no se presenta selectividad total por tener MicroLogic. La coordinación debe revisar los ajustes sin superar capacidad de conductores ni límites de UPS y receptores, y conservar detección de falla en inversor limitado. Las protecciones de motores ENV-19 requieren PTC autorizado y curva de arranque propia, por seleccionar, además del interruptor troncal. Las protecciones DC de cadenas, diferenciales finales, contactores, relés de subtensión y transferencia omnipolar permanecen partidas independientes, no suplidas por estos once conjuntos. \parencite{lote19-schneider-selectividad}.

\textbf{Concordancia normativa.} Se conserva íntegro el enunciado contractual que menciona NCh Elec. 2777. La Resolución Exenta 1.535 de 2009 identifica NCh2777.Of2003 como norma de seguridad de la información, anulada y reemplazada en esa resolución; no identifica una norma de puesta a tierra denominada NCh Elec. 2777. La oferta cumple el marco eléctrico vigente mediante Decreto 8 y pliegos RIC, especialmente 02/03/04/05/06/07/08 y 19, y registra la discrepancia literal para aclaración formal sin inventar equivalencia ni vigencia. Se entrega memoria/unifilar firmados por instalador competente, ensayos y declaración legal aplicable antes de puesta en servicio; son actuaciones futuras. ELE-19 individualiza los interruptores troncales; ELE-4 conserva RT-06.09 por completar por concordancia literal, coordinación/selectividad, protecciones finales y esquema de neutro real aún no resueltos \parencites[art.~2]{lote13-bcn-nch2777}{lote10-sec-ric}.

\subsubsection{Fronteras de potencia y actualización energética (ENE-4)}
\label{sec:sd4-energia-ene4}

ENE-4 conserva como envolvente de salida del camino superviviente 47,5076 kW iniciales y 48,4916 kW en crecimiento; la potencia aparente de crecimiento es 51,2335 kVA y su cociente conservador es 0,9465. Son máximos de diseño con margen, no consumo simultáneo de A+B ni factor de potencia medido. Los FP mínimos de compra por familia, los transitorios y las fases se mantienen en el balance. Una vez seleccionada la climatización de bancos, la nueva partida $P_b,S_b$ se incorpora con fuente y ubicación explícitas: $P_{sal,nuevo}=1,20(40,409679+P_{b,sal})$ y $S_{sal,nuevo}=1,20(42,694593+S_{b,sal})$ en crecimiento. Si antes de UPS, se incluye en demanda de fuente y continuidad propia; no se suma a salida. Deben cumplirse 60 kW/60 kVA, 20 kVA por fase y límites de corriente en todos los estados; no basta comparar una suma agregada con 60 kW.

Para presupuesto anual se conserva la frontera local del perfil de 8.760 horas: 8.798,40 kWh TI interior, 41.671,60 kWh del recinto y 68.250,87 kWh exportados. La razón de referencia es 4,7363, aproximadamente 4,74, condicionada a los parámetros de ventilación, humedad, UPS y clima ya declarados. La suma anual de energía importada al módulo en ese escenario es $41.671,60+68.250,87=109.922,47$ kWh; no se usa como numerador de PUE del recinto TI. Cloud, NOC, industria y ensayos externos conservan sus fronteras.

El efecto de bancos se incorpora por estado $j$, con horas $h_j$, electricidad adicional de ventiladores $F_{b,j}$, carga entálpica de reposición $Q_{b,j}$, calor adicional de carga/descarga $D_{b,j}$ y electricidad directa de control/tratamiento $H_{b,j}$. Ninguna carga entálpica se vuelve a sumar como sensible más latente. Si el tratamiento se alimenta después de las UPS y su electricidad/calor se retira por la misma frontera, se usa $k=1/0,960-1$ y el COP propio del escenario:
\[
\Delta C_j=\frac{F_{b,j}+H_{b,j}+[Q_{b,j}+D_{b,j}+F_{b,j}+H_{b,j}]/COP}{1-k/COP},\qquad
\Delta E_{rec}=\sum_j h_j(1+k)\Delta C_j.
\]
La pérdida UPS incremental $k\Delta C_j$ se cuenta una sola vez. Si hay refrigerador separado o tratamiento previo a UPS, se sustituye su rendimiento/camino y no se aplica esta expresión ciegamente. $D_b$ contiene solo calor adicional no incluido en pérdidas del modelo; la energía útil cargada que después abastece el mismo perfil no se vuelve a importar como pérdida anual.

Como sensibilidad reconstruible, no selección, se toma renovación de 300 m$^3$/h por banco durante 8.760 h, $Q_b=11,6774$ kW agregado desde la memoria entálpica y $F_b=0,600$ kW propio para ambos ventiladores. Con $D_b=H_b=0$ y COP propio 3, resulta $\Delta C=4,75856$ kW, $\Delta E_C=41.684,97$ kWh y $\Delta E_{rec}=43.421,84$ kWh. El PUE de ese escenario asciende a $(41.671,60+43.421,84)/8.798,40=9,6715$. Los 0,600 kW no son catálogo de un ventilador adquirido; omitir pérdidas adicionales de bancos hace esta sensibilidad incompleta para compra. El resultado muestra por qué el PUE 4,74 no puede trasladarse a renovación continua de bancos sin tratamiento elegido. La variante de extracción/tratamiento, horas y estados de recarga deben fijarse conjuntamente antes de emitir el balance final. ENE-4 desarrolla método, frontera y sensibilidad, pero conserva RT-06.11 por completar mientras esos consumos sigan indeterminados.

La recepción instala medición de importación por fuente, entrada/salida UPS, TI interior, cada sistema de clima/bancos y exportación. Los acumulados sincronizados conservan período y estados, y el ciclo anual compara reservas inicial/final equivalentes; los ensayos de bancada y la carga artificial de grupo se registran separadamente. El FP operacional se calcula con energía activa/reactiva y demanda aparente verdadera de la misma frontera, incluyendo distorsión; no se deduce del PUE ni de la suma de factores nominales. La actualización informa variaciones de parque, clima, pérdidas y uso sin presentar una medición o aceptación ya ejecutada \parencite[RT-06.11, p.~15]{pucv-transversal}.
